% Copyright 2023 Sycuricon Group
% Author: Phantom (phantom@zju.edu.cn)

\section[Appendix 1]{Appendix 1: Bit-Level Integer Coding Rules}\label{app1}

In \nameref{p4}, we will artificially restrict what programming constructs you can use to help you gain a better understanding of the bit-level, logic, and arithmetic operations of C. 
In answering these problems, your code must follow these rules:

\subsection*{Assumptions}

\begin{enumerate}
    \item Integers are represented in two's-complement form.
    \item Right shifts of signed data are performed arithmetically.
    \item Data type int is w bits long. For some of the problems, you will be given a specific value for \code{w}, but otherwise your code should work as long as w is a multiple of 8. You can use the expression \code{sizeof(int)<<3} to compute \code{w}.
\end{enumerate}

\subsection*{Forbidden}

\begin{enumerate}
    \item Conditionals (\code{if} or \code{?:}), loops, switch statements, function calls, and macro invocations. 
    \item Division, modulus, and multiplication.
    \item Relative comparison operators (\code{<}, \code{>}, \code{<=}, and \code{>=}).
\end{enumerate}

\subsection*{Allowed operations}

\begin{enumerate}
    \item All bit-level and logic operations.
    \item Left and right shifts, but only with shift amounts between \code{0} and \code{w-1}. Addition and subtraction.
    \item Equality (\code{==}) and inequality (\code{!=}) tests.
    \item Integer constants \code{INT\_MIN} and \code{INT\_MAX}. (TIP: \code{<limits.h>} should be included if you want to use them.)
    \item Casting between data types int and unsigned, either explicitly or implicitly.
\end{enumerate}

\noindent Even with these rules, you should try to make your code readable by choosing descriptive variable names and using comments to describe the logic behind your solutions.

\noindent As an example, the following code extracts the most significant byte from integer argument \code{x}:

\begin{lstlisting}[frame=tb, language=C++]
/* Get most significant byte from x */
int get_msb(int x) {
    /* Shift by w-8 */
    int shift_val = (sizeof(int)-1)<<3; /* Arithmetic shift */
    int xright = x >> shift_val;
    /* Zero all but LSB */
    return xright & 0xFF;
}
\end{lstlisting}





\section[Appendix 2]{Appendix 2: Bit-Level Floating-Point Coding Rules}\label{app2}

In \nameref{p5}, you will write code to implement floating-point functions, operating directly on bit-level representations of floating-point numbers.
Your code should exactly replicate the conventions for IEEE floating-point operations, including using round-to-even mode when rounding is required.

\noindent To this end, we define data type \code{float\_bits} to be equivalent to unsigned:

\begin{lstlisting}[frame=tb, language=C++]
/* Access bit-level representation floating-point number */
typedef unsigned float_bits;
\end{lstlisting}

\noindent Rather than using data type float in your code, you will use \code{float\_bits}. 
You may use both int and unsigned data types, including unsigned and integer constants and operations. 
You may not use any unions, structs, or arrays. 
Most significantly, you may not use any floating-point data types, operations, or constants. 
Instead, your code should perform the bit manipulations that implement the specified floating-point operations.

\noindent The following function illustrates the use of these coding rules. 
For argument \code{f} , it returns ±0 if \code{f} is denormalized (preserving the sign of \code{f}), and returns \code{f} otherwise.

\begin{lstlisting}[frame=tb, language=C++]
/* If f is denorm, return 0. Otherwise, return f */
float_bits float_denorm_zero(float_bits f) {
	/* Decompose bit representation into parts */ 
  unsigned sign = f>>31;
	unsigned exp = f>>23 & 0xFF;
	unsigned frac = f & 0x7FFFFF;
	if (exp == 0) {
    /* Denormalized. Set fraction to 0 */
		frac = 0; 
    }
	/* Reassemble bits */
	return (sign << 31) | (exp << 23) | frac; 
}
\end{lstlisting}
